% ============================================================
%  ch13  第五讲 §5.3 Langlands 对偶群（半无限轨道与 MV 理论）
%  原文行 2644–2929
% ============================================================

\subsection{Langlands 对偶群}\label{MV semiinfinite}

前面已见 $\Sat_G$ 连同交换约束 $c$ 与纤维函子 $\on{H}^*$ 构成中性 Tannaka 范畴（在命题 \ref{existence of comm} 之后成立，该命题将在下一小节证明）。令
\[\tilde G=\Aut^\otimes \on{H}^*\]
为 Tannaka 群。下一步是用 $G$ 显式描述这个仿射群概形。首先容易证明

\begin{lem}
$\Ql$ 上的仿射群概形 $\tilde G$ 连通且约化。
\end{lem}

\begin{proof}
首先注意 $\IC_{\mu_1}\star\IC_{\mu_2}$ 以 $\IC_{\mu_1+\mu_2}$ 为直和项，因为卷积映射 $\Gr_{\leq\mu_1}\tilde\times\Gr_{\leq \mu_2}\to \Gr_{\leq \mu_1+\mu_2}$ 双有理。由此可知 $\Sat_G$ 作为张量范畴由有限多个 $\IC_\mu$ 生成（例如由 $\mu$ 为基本余权者生成），且不存在只含有限多个 $\IC$ 的直和的张量子范畴（例如 $\IC_\mu$ 的 $n$ 重自卷积含 $\IC_{n\mu}$ 为直和项）。于是由 \cite[Proposition 2.20, Corollary 2.22]{DM}，$\tilde G$ 代数且连通。最后，因 $\Sat_G$ 半单，$\tilde G$ 约化（\cite[Proposition 2.23]{DM}）。
\end{proof}

回忆对一个连通约化群 $G$，可附上四元组 $(\xch,\xcoch,\Phi,\Phi^\vee)$，称为 $G$ 的\textbf{根数据}（参见 \S\ref{Intro:group}）。令人惊讶的事实是：代数闭域上的连通约化群在同构意义下由这个组合数据唯一确定。此外，若 $(\xch,\xcoch,\Phi,\Phi^\vee)$ 来自某个 $G$，则 $(\xcoch,\xch,\Phi^\vee,\Phi)$ 也来自某个 $\hat G$：即 $\xcoch$（相应地 $\xch$）是 $\hat G$ 的 Cartan 环面 $\hat T$ 的权格（相应地余权格），而 $\Phi^\vee\subset \xcoch$（相应地 $\Phi\subset\xch$）是相应的根系（相应地余根系）。综述见 \cite[\S~1,2]{Sp}。约化群 $\hat G$ 通常称为 $G$ 的 \textbf{Langlands 对偶群}。\zy{nLab 词条 Langlands dual group（\url{https://ncatlab.org/nlab/show/Langlands+dual+group}，2014-09-07 修订）目前只列参考文献、未给定义，故不引为定义依据；本条依原文所引综述 \cite{Sp}（A 级）。[来源：A 级 \cite{Sp}；nLab 词条现状如上。]}

注意从 $G$ 到 $\hat G$ 的过渡完全是组合性的，依赖约化群的分类定理。下表给出若干例子。

\begin{center}
\begin{tabular}{| c | c | c | c | c | c |}
\hline 
$G$ & $\GL_n$ & $\on{SL}_n$ & $\on{SO}_{2n+1}$ & $\on{SO}_{2n}$ & $E_8$  \\ \hline 
$\hat G$ & $\GL_n$ & $\on{PGL}_{n}$ & $\on{Sp}_{2n}$ & $\on{SO}_{2n}$ & $E_8$ \\ \hline 
\end{tabular}
\end{center}

\begin{thm}\label{birth of dual group}
Tannaka 群 $\tilde G$ 的根数据为 $(\xcoch,\xch,\Phi^\vee,\Phi)$，因此同构于 $\hat G$。
\end{thm}

\begin{rmk}
在某种意义上，本定理可看作 $G$ 的 Langlands 对偶群的无分类定义。
\end{rmk}

本小节余下部分解释该定理的证明思路与一些应用。在证明过程中将给 $\tilde G$ 添上更多结构（见命题 \ref{canonical pinning}）。注意 Richarz 给出了本定理的另一途径（参见 \cite[\S~3]{Ri}）。

首先容易验证

\begin{lem}
设 $G=T$ 是环面。则定理成立。
\end{lem}

\begin{proof}
事实上，$(\Gr_T)_{\on{red}}$ 是典范同构于 $\xcoch(T)$ 的离散点集。于是容易看出 $\Sat_T$ 等价于 $\xcoch(T)$-分次有限维 $\Ql$-向量空间范畴，而 $\on{H}^*$ 就是遗忘分次的函子。引理得证。
\end{proof}

为从这一特殊情形推出一般情形，模仿经典构造定义「范畴 Satake 变换」，即对称幺半函子
\[\on{CT}:\Sat_G\to \Sat_T,\]
其中 $T$ 是 $G$ 的抽象 Cartan。

沿用 \S\ref{Intro:group} 的记号。固定 Borel 子群 $B\subset G$，令 $T=B/U$ 为抽象 Cartan，其中 $U$ 是 $B$ 的幂幺根。

考虑下述诱导图
\begin{equation}\label{CT diagram}
\Gr_{T}\stackrel{q}{\leftarrow} \Gr_{B}\stackrel{i}{\to} \Gr_{G}.
\end{equation}
注意 $\Gr_{B}\to \Gr_{G}$ 在点层面是双射，但远非同构于归纳概形；事实上态射 $i$ 本身就很病态。不过 $(\Gr_B)_{\on{red}}$ 的每个连通分支在 $\Gr_G$ 中局部闭。由于只关心这些空间上的 $\ell$-进层，可安全地把涉及的仿射 Grassmann 簇换为其约化子归纳概形（但为简化记号，下文省略下标 $_{\on{red}}$）。于是如前所述，$\Gr_T\simeq \xcoch(T)$ 是离散点集。对 $\la\in\xcoch(T)$，定义
\[S_\la= i(q^{-1}(\la))\subset \Gr_G.\]
注意它是 $LU$ 在 $\Gr_G$ 上作用（经由嵌入 $LU\subset LG$）下过 $t^\la$ 的轨道。

回忆 $\xcoch(T)$ 上的偏序 $\leq$。首先有

\begin{prop}\label{geom of Sla}
并 $S_{\leq \la}:=\cup_{\la'\leq \la} S_{\la'}$ 闭，且 $S_\la\subset S_{\leq \la}$ 开稠密。特别地，$S_\la$ 在 $\Gr_G$ 中局部闭，而 $S_{\leq \la}$ 是其闭包 $\bar{S}_\la$。
\end{prop}

\begin{proof}
证明与命题 \ref{HP} 及命题 \ref{geom of Sch} 类似。先证 $\cup_{\la'\leq \la} S_{\la'}$ 闭。先设 $G_\der$ 单连通。则 $S_{\leq \la}$ 上的点有如下模解释：对 $G$ 的最高权表示 $V_\chi$，以 $\ell_\chi$ 记相应的最高权直线，则
\[S_{\leq \la}=\left\{(\mE,\beta)\in\Gr\mid \beta^{-1}(\ell_\chi)\subset t^{-(\chi,\la)}(V_{\chi,\mE}),\ \mbox{对一切 } V_\chi\right\}.\]
于是如命题 \ref{HP} 的论证，可把这些条件写成一组矩阵方程，从而定义闭子集。对一般 $G$，可通过 $z$-扩张化归到 $G_\der$ 单连通的情形。

剩下要证：若 $\al$ 是正余根，则 $S_{\la-\al}$ 在闭包 $\overline{S}_\la$ 中。这由 \eqref{curve C} 中曲线 $C_{\la,\al}$ 的构造推出。
\end{proof}

这些 $LU$-轨道与 Schubert 胞腔、对偶 Schubert 胞腔之间有下述关系。

\begin{lem}\label{intersection of semi and Sch}
若交 $S_\la \cap \Gr_\mu$ 非空，则 $t^\la\in \Gr_{\leq \mu}$；若交 $S_\la\cap \Gr^{\mu}$ 非空，则 $t^\la\in \Gr^{\geq \mu}$。
\end{lem}

\begin{proof}
以 $2\rho^\vee$ 记 $G$ 的正余根之和。选取嵌入 $T\subset B$。由复合
\[\bG_m\stackrel{2\rho^\vee}{\to} T\subset G\subset L^+G,\]
$L^+G$ 在 $\Gr_G$ 上的作用诱导 $\bG_m$ 在 $\Gr_G$ 上的作用。则 $\Gr_B$ 成为该作用的吸引子，而 $\Gr_T\subset \Gr_G$ 是不动点集，即
\[S_\la=\left\{x\in \Gr_G\mid \lim\limits_{s\to 0}2\rho^\vee(s)x=t^\la\right\}.\]
于是因 $\Gr_\mu$ 是 $2\rho^\vee(\bG_m)$-不变的，$S_\la\cap \Gr_\mu\neq\emptyset$ 仅当 $t^\la\in \overline{\Gr_\mu}=\Gr_{\leq \mu}$。同一论证给出第二个断言。
\end{proof}

\begin{cor}\label{boundary of semi}
集合 $S_{<\la}:=\cup_{\la'<\la} S_{\la'}$ 是 $\overline{S}_\la$ 与丰沛线丛 $\mO(1)$ 的某个超平面截影之交。
\end{cor}

当 $\cha\ k=0$ 时，Mirkovi\'c--Vilonen 用相应 Kac--Moody 代数（参见 \S\ref{Kac-Moody}）的基本表示 $L_1(\frakg_\Ga)$ 证明了这一点（参见 \cite[Proposition 3.1(b)]{MV}）。下面给出一个不依赖 Kac--Moody 理论（且对任意特征有效）的论证。

\begin{proof}
经平移 $t^{-\la}$，可设 $\la=0$。容易看出 $S_0\subset \Gr^0$。由引理 \ref{intersection of semi and Sch}，$S_{<0}\cap \Gr^0=\emptyset$。换言之，$S_{<0}=S_{\leq 0}\cap \Theta$，其中由定理 \ref{Pic of Gr} 回忆 $\Theta$ 是 $\Gr_G$（的中性连通分支）上对应于 $\mO(1)$ 的有效 Cartier 除子。证毕。
\end{proof}

现在可以叙述 Mirkovi\'c--Vilonen 关于半无限轨道的主要几何结果（参见 \cite[\S\ 3]{MV}）。

\begin{thm}\label{MV theory}~
\begin{enumerate}
\item 若 $t^\la\in \Gr_{\leq \mu}$，则 $S_\la\cap \Gr_\mu$ 非空，且 $S_\la \cap \Gr_{\leq \mu}$ 纯维数为 $(\rho,\la+\mu)$。

\item 设 $\mA\in\Sat_G$。则 $(q_!i^*\mA)_\la=\on{H}_c^*(S_\la, \mA)$ 集中在次数 $(2\rho,\la)$。此外有典范同构
\[\on{H}_c^{(2\rho,\la)}(S_\la, \IC_\mu)\otimes\Ql((\rho,\la+\mu))\cong \Ql[\on{Irr}(S_\la\cap\Gr_{\leq \mu})],\]
其中右端是 $S_\la\cap\Gr_{\leq \mu}$ 的不可约分支张成的 $\Ql$-空间。

\item 有函子的典范同构
\[\on{H}^*(-)\cong \on{CT}:=\bigoplus_\la \on{H}_c^*(S_\la,-):\Sat_G\to \on{Vect}_{\Ql}.\]
\end{enumerate}
\end{thm}

\begin{rmk}
Mirkovi\'c--Vilonen 对一般系数环 $\Lambda$ 建立了该定理（当 $k=\bC$）。在这种一般性下，(1) 中的 $\IC_\mu$ 须换为相应的标准或余标准对象。
\end{rmk}

下面说明该定理证明的基本思路，细节见 \cite[\S\ 3]{MV}。\cite[\S~2.2]{Z14} 基于 \cite{NP} 的某些想法给出了另一途径。这一替代途径对建立 $p$-进群的几何 Satake 是必要的，因为（目前）我们在那种情形下有相应的推论 \ref{boundary of semi}（那里既没有 Kac--Moody 理论，也没有 Cartier 除子 $\Theta$）。

首先，由推论 \ref{boundary of semi} 可知：若 $\al$ 是单余根，则 $\dim(S_{\la-\al}\cap\Gr_{\leq \mu})=\dim(S_\la\cap\Gr_{\leq \mu})-1$。于是由下述起点归纳即得 (1)：
\begin{equation}\label{extreme semi}
S_\la\cap \Gr_{\leq \mu}=S_\la\cap \Gr_\mu=\left\{\begin{array}{ll}t^{w_0(\mu)} & \la=w_0(\mu), \\
                               L^+U t^\mu\cong \bA^{(2\rho,\mu)} & \la=\mu, \end{array}\right.
\end{equation}
其中 $w_0$ 是 $G$ 的 Weyl 群的最长元。

至于 (2)，注意函子 $q_!i^*$ 恰是 Braden 的双曲局部化函子（参见 \cite{Br}）。更确切地说，设 $B^-$ 是满足 $B\cap B^-=T$ 的对偶 Borel，则 $\Gr_{B^-}$ 是 $2\rho^\vee(\bG_m)$-作用的排斥子，且有图
\[\Gr_T\stackrel{q^-}{\leftarrow}\Gr_{B^-}\stackrel{i^-}{\to}\Gr_G.\]
由 Braden 定理，对 $\mA\in\Sat_G$ 有
\[q_!i^*\mA=(q^-)_*(i^-)^!\mA.\]
此外，用 $\{S_\la\cap\Gr_\mu, \mu\in\xcoch(T)^+\}$ 对 $S_\la$ 分层，诱导一个 $E_1$-项为 $\on{H}_c^*(S_\la\cap\Gr_\mu,\mA)$、收敛到 $\on{H}_c^*(S_\la,\mA)$ 的谱序列。这些考虑连同 (1) 的维数公式即得 (2)。

为证 (3)，用半无限轨道 $\{S_\la, \la\in\xcoch(T)\}$ 对 $\Gr$ 分层，诱导一个 $E_1$-项为 $\on{H}_c^*(S_\la,-)$、收敛到 $\on{H}^*(-)$ 的谱序列。由 (ii) 中建立的理由，它在 $E_1$-项退化。于是 $\on{H}^*$ 上有自然滤过，其伴随分次为 $\oplus_\la \on{H}_c^*(S_\la,-)$。具体地说，这是以 $(\xcoch(T),\leq)$ 为指标、定义为
\[\on{Fil}_{\geq \mu}\on{H}^*(\mA)=\ker(\on{H}^*(\mA)\to \on{H}^*(S_{<\la},\mA))\]
的滤过，其中 $S_{<\la}=\bar{S}_\la-S_\la$。另一方面，用
\[\left\{S^-_\la:=i^-(q^-(\la)),\ \la\in\xcoch(T)\right\}\]
对 $\Gr$ 分层给出 $\on{H}^*$ 的另一个滤过
\[\on{Fil}'_{<\la}\on{H}^*(\mA)=\on{Im}(\on{H}^*_{S^-_{<\la}}(\mA)\to \on{H}^*(\mA)),\]
其中 $\on{H}^*_{S^-_{<\la}}(\mA)$ 记 $\mA$ 到 $S^-_{<\la}$ 的 $!$-限制的上同调。这两个滤过彼此互补（由 Braden 定理与 (2)），共同给出分解 $\on{H}^*=\oplus_\la \on{H}_c^*(S_\la,-)$。

\begin{rmk}
(i) 由此可知 $\dim (S_\la\cap \Gr_{\leq \mu})=(\rho,\la+\mu)$ 与 $\dim \Gr_{\leq \mu}-\dim (S_\la\cap \Gr_{\leq \mu})=(\rho,\mu-\la)$ 都随 $\mu$ 趋于无穷而趋于无穷。换言之，$S_\la$ 在 $\Gr_G$ 中既无穷维又无穷余维。另一方面它是 $LU$-轨道，因此 $S_\la$ 通常称为\textbf{半无限轨道}。

(ii) 交 $S_\la\cap S^-_\mu$ 有限维，且只依赖 $\la-\mu$，故通常记作 $Z_{\la-\mu}$，称为 \textbf{Zastava 空间}。注意 $Z_{\la-\mu}$ 非空当且仅当 $\la\geq \mu$。

(iii) \cite[Theorem 3.6]{MV} 还证明了 $\on{H}^*$ 的滤过与分解不依赖嵌入 $B\subset G$ 与 $T\subset B$ 的选取。另见推论 \ref{canonical TB}。
\end{rmk}

可把 $\on{CT}$ 看作函子 $\Sat_G\to \Sat_T$，把定理 \ref{MV theory} (iv) 中的同构看作
\begin{equation}\label{geom sat tran}
\on{H}^*\circ \on{CT}\cong \on{H}^*: \Sat_G\to \on{Vect}_\Ql.
\end{equation}

\begin{prop}
$\on{CT}$ 上有唯一幺半结构使同构 \eqref{geom sat tran} 幺半。
\end{prop}

$\on{CT}$ 上的幺半结构若存在则显然唯一。其存在性由 \cite[Proposition 6.4]{MV} 用卷积积的融合积解释证明。这里用等变上同调勾勒另一途径。

\begin{proof}
存在性如下。前已说明 $\on{H}^*$ 上有滤过 $\on{Fil}^*_{\geq \mu}$ 使 $\on{CT}=\on{gr}\on{H}^*$。只需证明命题 \ref{fiber functor} 给出的幺半结构与该滤过相容。于是由 Verdier 对偶论证知互补滤过也与幺半结构相容；再由上述讨论中 $\on{H}^*\cong \on{CT}$ 的构造即得本命题。

显然的困难是：$\on{H}^*$ 的幺半结构是用 $L^+G$-等变上同调构造的，而 $S_\la$ 并非 $L^+G$-等变。不过 $S_\la$ 在常值环面 $T\subset L^+T\subset L^+G$ 的作用下稳定。于是分层 $\{S_\la\}$ 也诱导一个 $E_1$-项为 $\on{H}_{T,c}^*(S_\la,-)$、收敛到 $\on{H}^*_{T}$ 的谱序列；由 (ii) 的理由它在 $E_1$-项退化，故可把滤过 $\on{Fil}_{\geq \mu}\on{H}^*$ 提升为 $\on{H}^*_{T}$ 的由自由 $R_{T}$-模构成的滤过
\[\on{Fil}_{\geq \mu}\on{H}_{T}^*(\mA)=\ker(\on{H}_{T}^*(\mA)\to \on{H}_{T}^*(S_{<\la},\mA)),\]
使伴随分次为 $\oplus_\la\on{H}_{T,c}(S_\la,\mA)$（特别地是自由 $R_{T}$-模）。此外，沿增广映射 $R_{T}\to\Ql$ 基变换后即恢复 $\on{H}^*$ 上的滤过。细节见 \cite[\S\ 2.2]{YZ}。

因 $\on{H}_{T}^*(\mA)\cong \on{H}^*_{L^+G}(\mA)\otimes_{R_{G}}R_{T}$，命题 \ref{monoidal functor} 中构造的 $\on{H}_{L^+G}^*$ 上的幺半结构经基变换给出 $\on{H}_{T}^*$ 上的幺半结构。因此只需证明典范同构
\begin{equation}\label{equiv monoidal}
\on{H}^*_{T}(\mA_1)\otimes_{R_{T}}\on{H}^*_{T}(\mA_2)\cong \on{H}^*_{T}(\mA_1\star\mA_2)
\end{equation}
与滤过相容；而这一点只需在 $\Spec R_{T}$ 的通用点处验证。令 $\on{H}_\eta=\on{H}^*_{T}\otimes_{R_{T}}Q$，其中 $Q$ 是 $R_{T}$ 的分式域。应用等变局部化定理（参见定理 \ref{equiv localization}）得
\[\on{H}_\eta(\mA)\cong \bigoplus_\nu \on{H}_\eta(\mA|_{t^\nu}),\quad \on{H}_\eta(S_{<\la},\mA)\cong\bigoplus_{\nu<\la}\on{H}_\eta(\mA|_{t^\nu}).\]
于是 $\on{Fil}_{\geq \la}\on{H}_\eta(\mA)=\oplus_{\nu\geq \la}\on{H}_\eta(\mA|_{t^\nu})$。此外，由 $\on{H}^*_{T}$ 上幺半结构的定义可知，在这些等同下 \eqref{equiv monoidal} 由典范同构
\[\begin{split}\bigoplus_{\nu_1+\nu_2=\nu} \on{H}_\eta(\mA_1|_{t^{\nu_1}})\otimes_Q \on{H}_\eta(\mA_2|_{t^{\nu_2}})\cong &\bigoplus_{\nu_1+\nu_2=\nu}\on{H}_\eta((\mA_1|_{t^{\nu_1}})\otimes (\mA_2|_{t^{\nu_2}})) \\
\cong &\quad \on{H}_\eta(\mA_1\star\mA_2|_{t^\nu})\end{split}\]
诱导，其中第一个同构是等变 K\"unneth 公式 \eqref{equiv Kunneth}，第二个仍是局部化定理。合起来即得命题。
\end{proof}

\begin{rmk}\label{Parabolic MV}
Mirkovi\'c--Vilonen 的理论可推广到抛物子群。设 $P\subset G$ 是抛物子群，$L_P$ 是其 Levi 商，则有诱导图
\[\Gr_L\stackrel{q}{\leftarrow}\Gr_P\stackrel{i}{\to}\Gr_G.\]
给定 $\mA\in\Sat_G$，可定义 $q_!i^*\mA$ 为 $\Gr_L$ 上的层复形。通过选取 $G$ 的适当单参数子群，也可把它看作双曲局部化函子，因此 $q_!i^*\mA$ 反常（差若干平移）。于是有 $\on{CT}=\on{CT}^G_T$ 的抛物版本
\[\on{CT}^G_L: \Sat_G\to \Sat_L,\quad \mA\mapsto q_!i^*\mA[\mbox{平移}],\]
满足 $\on{CT}^G_T=\on{CT}^L_T\circ\on{CT}^G_L$。
\end{rmk}

应用命题 \ref{existence of comm} 可见 $\on{CT}$ 事实上尊重对称幺半结构，从而是两个 Tannaka 范畴之间的张量函子，因此诱导同态
\[\hat{T}\cong \tilde T\to \tilde G.\]
这给出 $\tilde G$ 的一个子环面。

\begin{lem}
它是极大环面。
\end{lem}

\begin{proof}
事实上，令 $L_\mu=\on{H}^*(\IC_\mu)$ 为 $\tilde G$ 的不可约表示。则在定理 \ref{MV theory} (3) 给出的分解下，$\on{H}_c^{(2\rho,\la)}(S_\la,\IC_\mu)$ 是 $L_\mu$ 的 $\la$-权空间。由 \eqref{extreme semi}，$\dim \on{H}_c^{(2\rho,\mu)}(S_\mu,\IC_\mu)=1$，这说明 $L_\mu$ 的所有特征线性无关。因此 $\tilde G$ 的表示环是 $\tilde T$ 的表示环的子环，从而 $\tilde T$ 极大。
\end{proof}

\begin{rmk}
注意 $T$（定义为抽象 Cartan）不是 $G$ 的子群，但它的 Langlands 对偶 $\hat T$ 典范地是 $\tilde G$ 的子群。
\end{rmk}

继续定理的证明。令 $\ell_\mu=\on{H}_c^{(2\rho,\mu)}(S_\mu,\IC_\mu)$。注意自然映射（差一个纯量）
\[L_{\mu_1}\otimes L_{\mu_2}\to L_{\mu_1+\mu_2}\]
把 $\ell_{\mu_1}\otimes\ell_{\mu_2}$ 送到 $\ell_{\mu_1+\mu_2}$，因为 $\on{H}^*$ 上的幺半结构与 $\on{H}^*$ 上的滤过 $\on{Fil}$ 相容。于是由 Pl\"ucker 关系，族 $\{\ell_\mu\subset L_\mu, \mu\in\xcoch(T)^+\}$ 定义一个 Borel 子群 $\tilde B\subset \tilde G$，且 $\tilde T\subset \tilde B$。按定义，$\mu$ 是表示 $L_\mu$ 的（关于 $\tilde B$ 的）最高权。因此支配余权集 $\xcoch(T)^+$ 与 $\tilde T$ 关于 $\tilde B$ 的支配权集 $\xch(\tilde T)^+$ 一致。

回忆由 $\tilde G$ 的正根生成的半群 $\tilde Q^+\subset \xch(\tilde T)$ 的下述描述：$\la\in \tilde Q^+$ 当且仅当存在最高权表示 $L_\mu$ 使 $\mu-\la$ 作为 $L_\mu$ 的权出现。由定理 \ref{MV theory}，这等价于存在某个 $\mu$ 使 $t^{\mu-\la}\in \Gr_{\leq \mu}$，又等价于 $\la$ 是 $G$ 的正余根之和。因此由 $G$ 的正余根生成的半群 $(Q^\vee)^+\subset\xcoch(T)$ 与半群 $\tilde Q^+$ 一致，进而 $G$ 的单余根集与 $\tilde G$ 的单根集一致。为完成定理 \ref{birth of dual group} 的证明，只剩下

\begin{lem}
根数据 $(\xch,\xcoch,\Phi,\Phi^\vee)$ 由支配权半群 ${\xch}^+$ 与单根集 $\Delta\subset\Phi$ 唯一确定。
\end{lem}

\begin{proof}
事实上，${\xch}^+$ 确定 $\xch\otimes \bR$ 的支配 Weyl 房，从而确定 Weyl 群 $W$；于是 $\Phi=W(\Delta)$。细节留作练习。
\end{proof}

上述论证还给出

\begin{cor}\label{canonical TB}
超上同调函子 $\on{H}^*$ 上的分次对应于单参数子群
\[2\rho:\bG_m\to \tilde T\subset \tilde G,\]
其中余特征 $2\rho\in \xcoch(\tilde T)=\xch(T)$ 是 $\tilde G$ 的正余根之和。

特别地，$\tilde T$ 是 $2\rho(\bG_m)$ 在 $\tilde G$ 中的中心化子，而 $\tilde B$ 是包含 $\tilde T$ 且使 $2\rho$ 关于之支配的唯一 Borel 子群。因此 $(\tilde T,\tilde B)$ 不依赖 $T\subset B$ 与 $B\subset G$ 的选取。
\end{cor}

本小节最后给出几何 Satake 等价的若干应用。首先它给出 $\hat G$ 的表示的某些「典范基」。

\begin{thm}\label{MV basis}~
\begin{enumerate}
\item 经由定理 \ref{MV theory} (2)，$S_\la\cap \Gr_{\leq \mu}$ 的不可约分支给出 $L_\mu$ 的 $\la$-权空间的一组自然基。

\item 经由推论 \eqref{Satake fiber}，$\Gr_{\leq \mmu}\cap m^{-1}(t_\la)$ 中维数为 $(\rho,|\mmu|-\la)$ 的不可约分支给出 $\Hom_G(V_\la, V_{\mu_1}\otimes\cdots\otimes V_{\mu_n})$ 的一组自然基。
\end{enumerate}
\end{thm}

(1) 中的基通常称为 Mirkovi\'c--Vilonen 基，$S_\la \cap \Gr_{\leq \mu}$ 的不可约分支通常称为 MV 圈；(2) 中的基有时称为 Satake 基。

下面给出几何 Satake 等价对拓扑的几个漂亮应用。

\begin{ex}
(i) 有典范同构
\[\on{H}^*(\Gr(i,n))\cong \wedge^i \Big(\bigoplus_{j=0}^{n-1}\Ql[-2j](-j)\Big).\]

(ii) 设 $X_n\subset \bP^{n+1}$ 是光滑 $n$ 维二次超曲面，则
\[\on{H}^*(X_n)\cong\left\{\begin{array}{ll}
\bigoplus_{j=0}^n \Ql[-2j](-j) &  n \mbox{ 为奇数,} \\ 
\bigoplus_{j=0}^n \Ql[-2j](-j)\bigoplus \Ql[-n](-\frac{n}{2}) &  n \mbox{ 为偶数.}
\end{array}\right.\]
证明是把 $\Gr(i,n)$ 与 $X_n$ 分别等同于 $\on{PGL}_n$、$\on{SO}_{n+2}$ 的极小 Schubert 簇。
\end{ex}

回忆若 $k=\bC$，则 $\Gr\cong \Omega K$，从而是 $H$-空间。因此上同调环 $\on{H}^*(\Gr)$ 有自然的 $\bZ$-分次交换且余交换的 Hopf 代数结构。$k$ 一般时这仍成立：卷积映射 $m:\Gr\tilde\times \Gr\to \Gr$ 诱导等变上同调的余乘法
\[\on{H}^*_{L^+G}(\Gr)\to \on{H}^*_{L^+G}(\Gr)\otimes_{R_G}\on{H}^*_{L^+G}(\Gr).\]
引理 \ref{left=right} 中的类似陈述对 $\on{H}^*_{L^+G}(\Gr)$ 成立。因此可沿 $R_G\to\Ql$ 特殊化，得到通常上同调的余乘法。

由这类 Hopf 代数的分类，$\on{H}^*(\Gr)=U(\fraka)$（或当 $k=\bC$ 时 $\on{H}^*(\Omega K)=U(\fraka)$），其中 $\fraka$ 是 $\on{H}^*(\Gr_G)$ 的本原部分。因 $\on{H}^*(\Gr)$ 通过杯积作用于 $\on{H}^*(\mA)$，一般 Tannaka 形式化诱导映射 $\fraka\to \hat \frakg=\Lie \hat G$。

\begin{thm}\label{canonical pinning}
固定同构 $\Ql(1)\simeq \Ql$。
\begin{enumerate}
\item Chern 类 $c_1(\mO(1))\in \fraka$ 及其在 $\hat \frakg$ 中的像，是 $\hat \frakg$ 的主幂幺元 $N$。

\item 映射 $\fraka\to\hat\frakg$ 单射，并把 $\fraka$ 与 $N$ 在 $\hat\frakg$ 中的中心化子等同。

\item $(\tilde G,\tilde B,\tilde T, N)$ 是 $\tilde G$ 的典范带钉。
\end{enumerate}
\end{thm}

(1) 与 (2) 在 \cite[\S 4]{Gi1} 与 \cite{YZ} 中证明。这里只指出 $N$ 是主幂幺元是相交上同调的硬 Lefschetz 定理的推论（参见 \cite[Lemma 4.4.2]{Gi1}）。(3) 在 \cite{YZ} 中证明。该定理特别给出当 $k=\bC$ 时紧 Lie 群的基于环路群的描述。

\begin{rmk}
(i) 本定理有整系数版本，用 $\Gr$ 的同调表述（\cite[Theorem 1.1]{YZ}），证明依赖几何 Satake 等价的整系数版本 \cite{MV}。

(ii) 类似地可用 $\hat\frakg$ 的正则中心化子描述等变上同调 $\on{H}^*_{L^+G}(\Gr)$，见 \cite[Theorem 2.12]{BFM} 与 \cite[Proposition 6.6]{YZ}。这是研究 $L^+G$-等变导出范畴 $D_{L^+G}(\Gr)$（又称导出 Satake 范畴）的第一步，见 \cite{BF}。
\end{rmk}
